% Generated draft: requires mathematical and source review.
\documentclass[11pt]{article}
\usepackage[margin=25mm]{geometry}
\usepackage{fontspec}
\setmainfont{Latin Modern Roman}
\usepackage{amsmath,amssymb,mathtools,unicode-math}
\setmathfont{Latin Modern Math}
\usepackage{enumitem,needspace,xurl,hyperref}
\setlength{\parindent}{0pt}
\setlength{\parskip}{5pt}
\setlength{\emergencystretch}{3em}
\newcommand{\N}{\mathbb{N}}
\newcommand{\Z}{\mathbb{Z}}
\newcommand{\Q}{\mathbb{Q}}
\newcommand{\R}{\mathbb{R}}
\newcommand{\C}{\mathbb{C}}
\newcommand{\F}{\mathbb{F}}
\newcommand{\A}{\mathbb{A}}
\newcommand{\PP}{\mathbb{P}}
\newcommand{\E}{\mathbb{E}}
\newcommand{\eps}{\varepsilon}
\newcommand{\dd}{\,\mathrm{d}}
\begin{document}
\section*{UnsolvedMath ID: 30001213}
\textbf{Source problem code:} catalogue entry
\subsection{Definitions and mathematical statement}
The supplied question concerns an associative algebra \(A\) belonging to a flat family of algebras and a commutative Poisson algebra \(R\) arising as its semiclassical limit. The base field, the parameter space, the deformation data defining that limit, and the meaning of a suitably generic member are not specified. Accordingly, the statement below retains these unresolved qualifications rather than asserting a conjecture for all flat families.

For terminology, write \(k\) for the unspecified base field. A Poisson \(k\)-algebra is a commutative \(k\)-algebra \(R\) equipped with a \(k\)-bilinear Lie bracket \(\{-,-\}\) satisfying
\[
\{r,st\}=\{r,s\}t+s\{r,t\}\qquad(r,s,t\in R).
\]
An ideal \(I\subseteq R\) is a Poisson ideal if \(\{R,I\}\subseteq I\). Under the usual commutative Poisson convention, define
\[
\operatorname{PSpec}(R)=\{P\in\operatorname{Spec}(R):\{R,P\}\subseteq P\}.
\]
For an ideal \(J\subseteq R\), its Poisson core is
\[
\mathcal P(J)=\sum_{\substack{I\subseteq J\\I\text{ a Poisson ideal}}} I,
\]
the largest Poisson ideal contained in \(J\). Under the usual maximal-ideal-core convention, the Poisson-primitive spectrum is
\[
\operatorname{PPrim}(R)=\{\mathcal P(\mathfrak m):\mathfrak m\in\operatorname{MaxSpec}(R)\}.
\]
The supplied record does not explicitly specify these conventions or the hypotheses under which alternative definitions agree.

For \(A\), let \(\operatorname{Spec}(A)\) denote its proper two-sided prime ideals: a proper two-sided ideal \(P\) is prime if \(IJ\subseteq P\) implies \(I\subseteq P\) or \(J\subseteq P\) for all two-sided ideals \(I,J\) of \(A\). Let
\[
\operatorname{Prim}(A)=\{\operatorname{Ann}_A(M):M\text{ is a simple left }A\text{-module}\}.
\]
Use the usual Zariski, or hull-kernel, topology on each spectrum: closed sets consist of the ideals in that spectrum containing a prescribed ideal. These are the conventional topologies used to express the question; the supplied record does not define them explicitly.

\textbf{Qualified source question.} For \(A\) a suitably generic member of the specified flat family with semiclassical limit \(R\), do there exist homeomorphisms
\[
h_{\mathrm{prim}}:\operatorname{Prim}(A)\longrightarrow\operatorname{PPrim}(R),
\qquad
h_{\mathrm{prime}}:\operatorname{Spec}(A)\longrightarrow\operatorname{PSpec}(R)
\]
that preserve the appropriate dimensions and group actions?

Here the admissible families, genericity condition, dimension invariants, acting groups, and compatible actions remain to be specified. If a common group \(G\) and actions are intended, preservation of those actions means \(h(g\cdot I)=g\cdot h(I)\) for every \(g\in G\) and every ideal \(I\) in the relevant domain. No particular group or action is supplied. Thus the record determines a qualified correspondence question, but does not determine a fully quantified general conjecture with explicit hypotheses.
\subsection{Short English statement}
The question asks whether a suitably generic quantum algebra and its semiclassical Poisson limit have homeomorphic primitive spectra and homeomorphic prime spectra, with the correspondences preserving the intended dimensions and symmetries. The allowed deformations, genericity assumptions, dimensions, and group actions are unspecified, so a precise general scope requires further source clarification.
\subsection{Sources}
\begin{itemize}
\item Supplied local UnsolvedMath registry record, permanent ID \(30001213\), source problem code \texttt{catalogue entry}, entitled \emph{Semiclassical Correspondence of Primitive and Poisson Spectra}. Project registry locations: \texttt{data/unsolvedmath/problems.json} and supplemental \texttt{catalog/problems.json}. Attribution is to the Oberwolfach Reports Open Problems collection; no individual proposer is supplied. The record gives the proposed year as 2009.
\item Supplied source citation: \emph{Mini-Workshop: Non-Negativity is a Quantum Phenomenon} (2009), Oberwolfach Reports, pp.~767--800, DOI \texttt{10.4171/owr/2009/14}. Original report: \url{https://doi.org/10.4171/owr/2009/14}. The supplied background describes these reports as documenting mathematical workshops and their open-problem discussions.
\item Supplied literature source, labelled ``Goodearl semiclassical-limits survey'': \url{https://arxiv.org/abs/0812.1612}.
\item Supplied literature source, labelled ``Goodearl--Letzter primary special-case result'': \url{https://arxiv.org/abs/0708.1091}.
\item Supplied literature assessment, checked \texttt{2026-08-21}: the record labels the problem partially solved and reports that the semiclassical primitive-spectrum conjecture is proved for several families, including quantum affine spaces and related toric cases, while the report and subsequent survey retain a general conjectural relationship. This is supplied dated curation derived from the OWR collection, not an independently verified literature assessment. The background explicitly requests independent verification of time-sensitive status and cited claims.
\item The supplied statement status is \texttt{corrected\_verified}; its accompanying verification describes meaning-preserving semantic, typographic, and LaTeX cleanup with no unsupported source-dependent additions. This is a supplied editorial designation, not independent mathematical verification of the formulation or conjecture.
\end{itemize}
\end{document}
